\documentclass[a4paper,12pt]{article}
\usepackage{amsmath}
\usepackage[russian]{babel}
\usepackage[utf8]{inputenc}
\usepackage[T2A]{fontenc}
\usepackage{geometry}
\usepackage{listings}
\usepackage{xcolor}
\usepackage{graphicx}
\usepackage{color}


\geometry{margin=1.5cm}
% Пример темы в стиле тёмной IDE
\definecolor{intj-bg}{HTML}{2B2B2B}
\definecolor{intj-key}{HTML}{CC7832}
\definecolor{intj-str}{HTML}{6A8759}
\definecolor{intj-comm}{HTML}{808080}
\definecolor{intj-id}{HTML}{A9B7C6}
\definecolor{intj-num}{HTML}{6897BB}

\lstdefinestyle{intellij}{
    language=Python,
    backgroundcolor=\color{intj-bg},
    basicstyle=\ttfamily\small\color{intj-id},
    keywordstyle=\color{intj-key}\bfseries,
    stringstyle=\color{intj-str},
    commentstyle=\itshape\color{intj-comm},
    numberstyle=\color{intj-comm}\tiny,
    stepnumber=1,
    frame=single,
    rulecolor=\color{intj-comm},
    showstringspaces=false,
    breaklines=true,
    tabsize=4
}






\begin{document}

% ---------- Титульный лист ----------
    \begin{titlepage}
        \begin{center}
            \textbf{Федеральное государственное автономное образовательное учреждение высшего образования «Национальный исследовательский университет ИТМО} \\
            \vspace{1cm}

            \textbf{Факультет программной инженерии и компьютерной техники} \\
            \vspace{1cm}

            \Large{\textbf{ОТЧЁТ}} \\
            \Large{\textbf{по задачам: I-L}} \\
            \large{по дисциплине «Алгоритмы и структуры данных»} \\
            \vspace{2cm}
        \end{center}

        \begin{flushright}
            \textbf{Выполнил:} \\
            Ануфриев Андрей Сергеевич \\
            P3219 \\
        \end{flushright}

        \vfill
        \begin{center}
            Санкт-Петербург 2026
        \end{center}
    \end{titlepage}

% ---------- Содержание ----------
    \tableofcontents
    \newpage

% ========== ЗАДАНИЕ 1 ==========


    \section{Задание I. Машинки}

    \lstinputlisting[style=intellij, label=lst:taskI]{../tasks/I.cpp}

    \subsection*{Суть решения}
    \textbf{Ключевая идея:} оптимальная кеш-стратегия - когда пол переполнен, убрать машинку, которая будет запрошена позже всех (или вообще больше не понадобится). Для этого предварительно считаем для каждой позиции, когда машинка потребуется в следующий раз.\\
    \textbf{Основной алгоритм:}
    \begin{enumerate}
        \item \textbf{Предвычисление next\_pos}: проходим последовательность запросов в обратном порядке. Для каждой позиции i считаем next\_pos[i] - позицию следующего вхождения этой машинки в будущем. last\_occurrence хранит последнее (т.е. ближайшее в будущем) вхождение каждой машинки.
        \item \textbf{Поддержание машинок на полу}: используем set<pair<int, int>> на\_полу с ключом (текущее\_время\_следующего\_вхождения, номер\_машинки). Set автоматически держит машинку с наибольшим временем следующего вхождения.
        \item \textbf{Обработка каждого запроса i}:
        \begin{itemize}
            \item Если машинка уже на полу: удаляем её из set со старым ключом и добавляем с новым (время следующего вхождения из next\_pos[i]).
            \item Если машинки нет на полу:
            \begin{itemize}
                \item Увеличиваем счётчик операций.
                \item Если пол переполнен, удаляем из конца set машинку, которая будет запрошена позже всех.
                \item Добавляем новую машинку с её временем следующего вхождения из next\_pos[i].
            \end{itemize}
        \end{itemize}
    \end{enumerate}
    По сути я просто реализовал Belady's optimal algorithm для данной задачи.\\
    \textbf{Время: O(P log K)} - P запросов, для каждого одна операция с set размером до K.\\
    \textbf{Память: O(N + P)} - массивы next\_pos и on\_floor, последовательность запросов.


% ========== ЗАДАНИЕ 2 ==========
    \newpage


    \section{Задание J. Гоблины и очереди}

    \lstinputlisting[style=intellij, label=lst:taskJ]{../tasks/J.cpp}

    \subsection*{Суть решения}
    \textbf{Ключевая идея:} разделить очередь на две половины с помощью двух деков, поддерживая инвариант left.size() >= right.size(). Это позволяет вставлять гоблинов в условную середину за O(1) за счёт операций переноса элементов между деками.\\
    \textbf{Основной алгоритм:}
    \begin{enumerate}
        \item \textbf{Инвариант структуры}: поддерживаем два дека (left и right). Начало очереди - front левого дека.
        \item \textbf{Операция +x}:
        \begin{itemize}
            \item Добавляем x в конец right.
            \item Проверяем баланс: если right стал больше left, то переносим элемент из front right в back left, чтобы восстановить инвариант.
        \end{itemize}
        \item \textbf{Операция *x}:
        \begin{itemize}
            \item Если left.size() > right.size() (нечётная длина очереди), добавляем x в front right - это ровно за центром.
            \item Если left.size() == right.size() (чётная длина), добавляем x в back left - это сразу в центр.
        \end{itemize}
        \item \textbf{Операция - }:
        \begin{itemize}
            \item Выводим front left (первого гоблина) и удаляем его.
            \item Проверяем баланс.
        \end{itemize}
    \end{enumerate}
    \textbf{Время: O(N)} - N операций, каждая за O(1).\\
    \textbf{Память: O(N)} - два дека суммарно содержат N элементов.


% ========== ЗАДАНИЕ 3 ==========
    \newpage


    \section{Задание K. Менеджер памяти-1}

    \lstinputlisting[style=intellij, label=lst:taskK]{../tasks/K.cpp}

    \subsection*{Суть решения}
    \textbf{Ключевая идея:} представить память как множество свободных блоков (размер, начало). При выделении найти первый подходящий блок, при освобождении слить освобождённый блок с соседними свободными блоками.\\
    \textbf{Основной алгоритм:}
    \begin{enumerate}
        \item \textbf{Инициализация}: начально вся память свободна - добавляем в set a блок (размер N, начало 0). Так же используем maps b и c для быстрого поиска соседей по началу и концу блока.
        \item \textbf{Запрос выделения K ячеек}:
        \begin{itemize}
            \item Получаем наибольший свободный блок из set rbegin().
            \item Если его размер z равен K, выделяем весь блок.
            \item Если z > K, выделяем первые K ячеек начиная с позиции y, а остаток (z-K ячеек начиная с y+K) остаётся свободным.
            \item Если z < K, запрос отклоняется - выводим -1.
        \end{itemize}
        \item \textbf{Запрос освобождения памяти}:
        \begin{itemize}
            \item Если запрос T был отклонен, игнорируем.
            \begin{itemize}
                \item Проверяем наличие блока слева, заканчивающегося в позиции v.second - 1. Если есть, объединяем его с текущим.
                \item Проверяем наличие блока справа, начинающегося в позиции v.second + v.first. Если есть, объединяем его с текущим.
            \end{itemize}
            \item Добавляем объединённый блок обратно.
        \end{itemize}
    \end{enumerate}
    Тут у меня было много попыток оптимизировать время. В итоге нашёл эту задачу в интернете и посмотрел на неё разбор. Написал примерно такое-же решение.\\
    \textbf{Время: O(M log M)} - M запросов, для каждого одна-две операции с set.\\
    \textbf{Память: O(M)} - хранение информации о выделенных блоках и свободных сегментах.


% ========== ЗАДАНИЕ 4 ==========
    \newpage


    \section{Задание L. Минимум на отрезке}

    \lstinputlisting[style=intellij, label=lst:taskL]{../tasks/L.cpp}

    \subsection*{Суть решения}
    \textbf{Ключевая идея:} использовать двустороннюю очередь для эффективного поддержания минимума в скользящем окне. Дек хранит индексы элементов текущего окна, упорядоченные по значениям в возрастающем порядке.\\
    \textbf{Основной алгоритм:}
    \begin{enumerate}
        \item \textbf{Инвариант дека}: дек содержит индексы элементов, находящихся в текущем окне, в порядке возрастания их значений. Минимум текущего окна всегда находится в front дека.
        \item \textbf{Удаление устаревших индексов}:
        \begin{itemize}
            \item Перед обработкой элемента i проверяем, не устарел ли индекс в front дека.
            \item Условие: dq.front() <= i - K означает, что индекс вышел за левую границу окна [i-K+1, i].
            \item Если устарел, удаляем его из front дека.
        \end{itemize}
        \item \textbf{Добавление нового элемента}: добавляем индекс i в конец дека. Теперь дек снова упорядочен по возрастанию значений.
    \end{enumerate}
    Изначально хотел сделать 2 указателя на массиве, но это было бы O(N*K) в худшем случае. Поэтому взял просто двустороннюю очередь.\\
    \textbf{Время: O(N)} - N элементов, каждый обрабатывается константное число раз.\\
    \textbf{Память: O(K)} - дек содержит не более K элементов (размер окна).

\end{document}